\section{Economic Significance}
\label{sec:economic}
We translate our statistical forecasting accuracy into economic utility via a volatility-targeting strategy. The portfolio scales its leverage inversely to the 1-day ahead volatility forecast to target a constant risk budget (2\% daily volatility), constrained by a maximum leverage of 5x. We assume a Quarter-Kelly position sizing framework and realistic transaction costs of 5 basis points (bps) per trade.

Table \ref{tab:backtest} reports the out-of-sample backtest metrics. While Neural-HAR provides statistically superior volatility forecasts (particularly for SPX), this does \textit{not} translate to a higher Sharpe ratio than the baseline HAR model. The sentiment-aware Neural-HAR model exhibits significantly higher turnover (e.g., 0.365 for SPX compared to 0.093 for HAR) as it aggressively reacts to non-linear signals. Net of the 5 bps slippage, this turnover erosion degrades the risk-adjusted return.

However, from a pure risk-management perspective, the Neural-HAR model is economically valuable. By providing more accurate volatility forecasts, the Neural-HAR strategy produces lower realized annualized volatility (closer to the target mandate) and smaller maximum drawdowns for both BTC ($-59.5\%$ vs HAR's $-65.4\%$) and SPX ($-35.2\%$ vs HAR's $-42.3\%$). Thus, Neural-HAR offers better risk control, albeit at the cost of total return due to implementation friction.

\begin{table}[htbp]
\centering
\caption{Economic Significance: Volatility-Targeting Backtest (5\,bps Slippage, Quarter-Kelly, $\leq 5{\times}$ Leverage, $N\approx275$--$287$)}
\label{tab:backtest}
\resizebox{\textwidth}{!}{%
\begin{tabular}{llrrrrrrrr}
\toprule
Asset & Model & Sharpe & Ann\textbackslash \_Return & Ann\textbackslash \_Vol & MaxDD & Calmar & HitRate & Turnover & N \\
\midrule
BTC & HAR & -1.994300 & -78.550000 & 39.390000 & -65.400000 & -1.201000 & 0.461800 & 0.073700 & 275.000000 \\
BTC & HAR-S & -2.031200 & -85.140000 & 41.920000 & -71.080000 & -1.197800 & 0.461800 & 0.057500 & 275.000000 \\
BTC & Neural-HAR & -2.095700 & -70.940000 & 33.850000 & -59.510000 & -1.192200 & 0.461800 & 0.052700 & 275.000000 \\
SPX & HAR & 1.446200 & 83.370000 & 57.640000 & -42.260000 & 1.972700 & 0.564500 & 0.093500 & 287.000000 \\
SPX & HAR-S & 1.161500 & 64.710000 & 55.720000 & -50.500000 & 1.281500 & 0.550500 & 0.579600 & 287.000000 \\
SPX & Neural-HAR & 1.201100 & 48.770000 & 40.610000 & -35.240000 & 1.384100 & 0.554000 & 0.365500 & 287.000000 \\
NIFTY & HAR & -0.310700 & -16.480000 & 53.050000 & -67.280000 & -0.245000 & 0.510900 & 0.050700 & 276.000000 \\
NIFTY & HAR-S & -0.383100 & -22.570000 & 58.910000 & -77.110000 & -0.292700 & 0.510900 & 0.165200 & 276.000000 \\
NIFTY & Neural-HAR & -0.572400 & -25.850000 & 45.160000 & -64.750000 & -0.399200 & 0.510900 & 0.484400 & 276.000000 \\
\bottomrule
\end{tabular}
}
\end{table}

